% !TeX spellcheck = ru_RU
% !TEX root = vkr.tex

\section{Описание решения}

\subsection{Используемые инструменты}

Традиционно для создания eDSL используются гомоиконичные языки (см.~\cite{HomoiconicLanguages}), например Lisp (и его производные) или TCL.
В данном случае их использование было затруднено, в силу необходимости работы с документами Office Open XML.

Кроме гомоиконичных языков для создания eDSL также часто используют функциональные языки программирования.
Поэтому в силу наличия уже существующего инструментария для обработки учебных планов\footnote{Инструменты для упрощения бумажной работы Юрия Викторовича Литвинова: \url{https://github.com/yurii-litvinov/TeachingScripts} (дата обращения: 11.02.2026).}, а также \enquote{простоты} работы с Office Open XML, был выбран \fsharp{}.
Реализация eDSL на \fsharp{} активно полагается на инструмент языка называемый Computation Expressions (далее~--- CE).

Также для экспорта учебного плана в Office Open XML Workbook используется дружественная к пользователю библиотека-обёртка над Open XML SDK\footnote{Репозиторий: \url{https://github.com/dotnet/Open-XML-SDK} (дата обращения: 11.02.2026).}~--- ClosedXML\footnote{Репозиторий: \url{https://github.com/ClosedXML/ClosedXML} (дата обращения: 11.02.2026).}.

\subsection{Подробности реализации}

В рамках работы сначала была реализована модель данных в соответствии с описанием из раздела~\ref{sec:relatedworks}.

\begin{listing}
      \begin{minted}[breaklines]{fsharp}
simpleBlock {
    workload 2
    competencies [ "ОПК-1"; "ПКА-1" ]

    disciplines
        [ discipline {
                number 002188
                name "Теория вероятностей и математическая статистика"
                englishName "Probability Theory and Mathematical Statistics"
                realization "осн курс"
                trajectory "тр 3 г"
                assessmentForms [ Credit ]

                classroomWork (
                    // Коллизия имён!
                    ClassroomWorkBuilder() {
                        lectures 30
                        practicalClasses 12
                        controlWorks 2
                        intermediateAssessment 2
                    }
                )

                independentWork (
                    // Коллизия имён!
                    IndependentWorkBuilder() {
                        usingMaterials 18
                        intermediateAssessment 8
                    }
                )

                interactiveHours 12
            } ]
}
      \end{minted}
      \caption{Описание обязательного предмета в первой версии eDSL}
      \label{lst:eDSL1}
\end{listing}

Затем в несколько итераций был разработан сам eDSL (на примере~\cite{CreatingDSLsUsing2022}).
В первой своей версии (см. листинг~\ref{lst:eDSL1}) язык почти дословно повторял описание модель данных, однако, во-первых, оказался сильно многословен, во-вторых, случились коллизии имен CE.
Во второй версии (см. листинг~\ref{lst:eDSL2}) была исправлена коллизия имен CE, а также созданы дополнительные CE, упрощающие жизнь пользователя (например, \texttt{simpleDiscipline} для описания обязательных предметов).

\begin{listing}
      \begin{minted}[breaklines]{fsharp}
simpleDiscipline {
    workload 2
    competencies [ "ОПК-1"; "ПКА-1" ]
    number 002188
    name "Теория вероятностей и математическая статистика"
    englishName "Probability Theory and Mathematical Statistics"
    realization "осн курс"
    trajectory "тр 3 г"
    assessmentForms [ credit ]

    classroomWork {
        lectures 30
        practicalClasses 12
        controlWorks 2
        intermediateAssessment 2
    }

    independentWork {
        usingMaterials 18
        intermediateAssessment 8
    }

    interactiveHours 12
}
      \end{minted}
      \caption{Описание обязательного предмета во второй версии eDSL}
      \label{lst:eDSL2}
\end{listing}

И в конце был описан экспорт в Office Open XML Workbook.
